% ------------------------------------------------------------
% English working translation layer.
% Source file: papers_30_59/58_Structures_de_Hodge_mixtes_reelles/58_Hodge_Mixtes_Reelles.tex
% Modern title: Real mixed Hodge structures
% ------------------------------------------------------------

% Structures de Hodge mixtes reelles
% Pierre Deligne
% Published in: Proceedings of Symposia in Pure Mathematics, Volume 55 (1994), Part 1, pp. 509-514
% This is a faithful LaTeX typesetting of the original paper
\documentclass[12pt,a4paper]{article}

% Encoding and fonts
\usepackage{fontspec}
\usepackage[english]{babel}

% Math packages
\usepackage{amsmath, amssymb, amsthm, amscd}
\usepackage{mathtools}
\usepackage{mathrsfs}

% Layout
\usepackage[margin=2.5cm]{geometry}
\usepackage{enumitem}
\usepackage{graphicx}
\usepackage{xcolor}

% Hyperlinks
\usepackage{hyperref}

% Theorem environments
\theoremstyle{plain}
\newtheorem{lemme}{Lemma}[section]
\newtheorem{proposition}[lemme]{Proposition}

\theoremstyle{definition}
\newtheorem{construction}[lemme]{Construction}
\newtheorem{remarque}[lemme]{Remark}

% Named operators
\DeclareMathOperator{\Gr}{Gr}
\DeclareMathOperator{\Aut}{Aut}
\DeclareMathOperator{\End}{End}
\DeclareMathOperator{\pr}{pr}

% Shortcuts
\newcommand{\A}{\mathcal{A}}
\newcommand{\Gm}{\mathbb{G}_\mathrm{m}}
\newcommand{\C}{\mathbb{C}}
\newcommand{\R}{\mathbb{R}}
\newcommand{\Z}{\mathbb{Z}}

\begin{document}
\begin{center}\small\emph{Literal English translation working draft. Formulae, numbering, and TeX structure are preserved; this is not yet a modern recast.}\end{center}
\vspace{0.5em}


% Title page header
\begin{center}
\small\textsc{Proceedings of Symposia in Pure Mathematics}\\
\small Volume 55 (1994), Part 1
\end{center}

\vspace{1cm}

\begin{center}
\Large\textbf{Mixed Hodge structures r\'{e}they}
\end{center}

\begin{flushright}
\textsc{Pierre Deligne}
\end{flushright}

\vspace{0.5cm}

The mixed Hodge structures r\'{e}they (\S2) forment a $\otimes$-cat\'{e}gorie ab\'{e}lienne
rigide and the functor ``space vector r\'{e}el under-jacent'' is a functor
fiber. The functor $\Gr^W$ is a other. We we proposons of calculer this that
given the th\'{e}orie g\'{e}n\'{e}rale of the cat\'{e}gories tannakiennes in this case particulier.

\section{Trois filtrations oppos\'{e}es}

\subsection{}

This num\'{e}ro is a variation on the num\'{e}ro 1.2 of P.~Deligne, \textit{Th\'{e}or\`{e}me
of Hodge~II}, Publ.~Math.~IHES \textbf{40}, pp.~5--58. Let $\A$ a cat\'{e}gorie ab\'{e}li-\-enne,
and $A$ a object of $\A$ equipped with trois filtrations $W$, $F$, $\bar{F}$ oppos\'{e}es.
``Filtration'' signifie ici ``filtration finite'' (loc.~cit.~1.1.4), the filtration $W$ is
suppos\'{e}e croissante, and the filtrations $F$, $\bar{F}$ d\'{e}croissantes. The condition
of opposition s'\'{e}crit $\Gr_F^p \Gr_{\bar{F}}^q \Gr_n^W A = 0$ for $n \neq p+q$
(loc.~cit.~1.2.7, 1.2.13).

For every object bigradu\'{e} $B = \bigoplus B^{p,q}$, we shall denote $W$, $F_W$, and $\bar{F}_W$
the trois filtrations oppos\'{e}es:
\begin{align*}
W_n &= \bigoplus_{p+q \leq n} B^{p,q}, \\
F_W^i &= \bigoplus_{p \geq i} B^{p,q}, \\
\bar{F}_W^i &= \bigoplus_{q \geq i} B^{p,q}.
\end{align*}
With this notation, that trois filtrations $W$, $F$, and $\bar{F}$ of $A$ let oppos\'{e}es
signifie again (loc.~cit.~1.2.6) that each $\Gr_n^W A$ admet a (unique)
d\'{e}composition
\begin{equation*}
\Gr_n^W A = \bigoplus_{p+q=n} A_W^{p,q}
\end{equation*}
such that the filtrations of $\Gr_n^W A$ induites by $F$ and $\bar{F}$ let donn\'{e}es by
\begin{align*}
F^i \Gr_n^W A &= \bigoplus_{\substack{p+q=n \\ p \geq i}} A_W^{p,q}, \\
\bar{F}^i \Gr_n^W A &= \bigoplus_{\substack{p+q=n \\ q \geq i}} A_W^{p,q}.
\end{align*}
En of other termes, $F$ and $\bar{F}$ induisent on $\Gr_n^W A$ the filtrations $F_W$ and $\bar{F}_W$.

Posons
\begin{equation*}
A_F^{p,q} := \bigl(W_{p+q} \cap F^p\bigr) \cap \Bigl( \bigl(W_{p+q} \cap \bar{F}^q\bigr) + \sum_{i>0} (W_{p+q-i} \cap \bar{F}^{q+1-i}) \Bigr)
\end{equation*}
and, \'{e}changeant the r\^{o}the of $F$ and $\bar{F}$ and of $p$ and $q$,
\begin{equation*}
A_{\bar{F}}^{p,q} := \bigl(W_{p+q} \cap \bar{F}^q\bigr) \cap \Bigl( \bigl(W_{p+q} \cap F^p\bigr) + \sum_{i>0} (W_{p+q-i} \cap F^{p+1-i}) \Bigr).
\end{equation*}

For $p+q=n$, the projection of $W_n$ on $\Gr_n^W A$ induces a isomorphism
of $A_F^{p,q}$ with $A_W^{p,q}$: loc.~cit.~1.2.8, restreint to the cas particulier 1.2.11 (o\`{u} $A_F^{p,q}$
is not\'{e} $A_0^{p,q}$; noter that in the formule d\'{e}finissant $A_F^{p,q}$, one can remplacer
the sum on $i > 0$ by the m\^{e}me sum on $i \geq 2$). One en d\'{e}duit (loc.
cit.) that the $A_F^{p,q}$ (resp.~$A_{\bar{F}}^{p,q}$) forment a bigrading of $A$, and that the
sum $a_F$ (resp.~$a_{\bar{F}}$) of the isomorphisms $A_F^{p,q} \xrightarrow{\sim} A_W^{p,q}$
(resp.~$A_{\bar{F}}^{p,q} \xrightarrow{\sim} A_W^{p,q}$)
is a isomorphism bifiltr\'{e} of $(A; W, F)$ with $(\Gr^W\!A; W, F_W)$ (resp.~of
$(A; W, \bar{F})$ with $(\Gr^W\!A; W, \bar{F}_W)$).

\begin{lemma}
the automorphism $d = a_{\bar{F}} a_F^{-1}$ of $\Gr^W\!A$ v\'{e}rifie
\begin{equation}\label{eq:1.1.1}
(d-1)(A_W^{p,q}) \subset \bigoplus_{\substack{r < p \\ s < q}} A_W^{r,s}.\tag{1.1.1}
\end{equation}
\end{lemma}

By passage to the gradu\'{e}, the isomorphisms $a_F$ and $a_{\bar{F}}$ induisent the identit\'{e}
of $\Gr^W\!A$. one has therefore $\Gr^W(d) = 1$.

For each part $I$ of $\Z^2$, and for $B$ a object bigradu\'{e}, posons $B_I =
\bigoplus_{(p,q) \in I} B^{p,q}$. Posons
\begin{align*}
I(p,q) &= \{(r,s) \mid r+s \leq p+q \text{ et } (r,s) = (p,q) \text{ ou } r < p\}, \\
J(p,q) &= \{(r,s) \mid r+s \leq p+q \text{ et } (r,s) = (p,q) \text{ ou } s < q\}, \\
K(p,q) &= \{(r,s) \mid r+s \leq p+q \text{ et } r \geq p\}.
\end{align*}

\begin{center}
\begin{tabular}{ccc}
\begin{tabular}{c}
$I(p,q)$: \\[6pt]
$J(p,q)$: \\[6pt]
$K(p,q)$: \text{ entour\'{e}}
\end{tabular}
&
\begin{tabular}{c@{\ }c@{\ }c@{\ }c@{\ }c}
& & & $\times$ & \\
& & $\times$ & $\times$ & \\
& $\times$ & $\times$ & $\circ$ & $(p,q)$ \\
$\circ$ & & & $\circ$ & \\
$\circ$ & $\circ$ & & $\circ$ &
\end{tabular}
&
\qquad
\begin{tabular}{c@{\ }c@{\ }c@{\ }c@{\ }c}
& & & $\circ$ & \\
& & $\circ$ & $\circ$ & \\
& $\circ$ & $\circ$ & $\times$ & $(p,q)$ \\
$\times$ & & & $\circ$ & \\
$\times$ & $\times$ & & $\circ$ &
\end{tabular}
\end{tabular}
\end{center}

In the d\'{e}finition of $A_F^{p,q}$ as intersection of deux under-objects of $A$, the
prime under-object is $A_{F}^{I(p,q)}$ and the second is $A_{\bar{F}}^{J(p,q)}$. That r\'{e}sulte of this that
$a_F$ (resp.~$a_{\bar{F}}$) is a isomorphism bifiltr\'{e}. one has therefore $A_F^{p,q} \subset A_F^{I(p,q)}$,
of o\`{u}
$d A_F^{p,q} = a_{\bar{F}} a_F^{-1} A_F^{p,q} \subset A_{\bar{F}}^{I(p,q)}$, and $d A_F^{p,q} \subset A_{\bar{F}}^{J(p,q)}$. By this that $\Gr^W(d) = 1$, one has \'{e}galit\'{e}: $d A_F^{I(p,q)} = A_{\bar{F}}^{I(p,q)}$. \'{E}changeant $F$ and $\bar{F}$, one trouve of m\^{e}me that
$d^{-1} A_{\bar{F}}^{J(p,q)} = A_F^{J(p,q)}$, therefore that $d A_F^{J(p,q)} = A_{\bar{F}}^{J(p,q)}$, and
\begin{equation*}
d A_W^{p,q} = d A_W^{I(p,q)} \cap d A_W^{J(p,q)} = A_W^{I(p,q) \cap J(p,q)}
= A_W^{K(p,q)} = \bigoplus_{\substack{r < p \\ s < q}} A_W^{r,s}.
\end{equation*}

One conclut en utilisant that $\Gr^W(d) = 1$:

\begin{proposition}\label{prop:1.2}
The functor $A \mapsto \Gr^W\!A$ is a \'{e}quivalence of the cat\'{e}gorie of the objects of $\A$, equipped with trois filtrations oppos\'{e}es $W$, $F$, $\bar{F}$, with the
cat\'{e}gorie of the objects of $\A$, equipped with a bigrading (under-entendu: finite) and
of a automorphism $d$ v\'{e}rifiant \eqref{eq:1.1.1}.
\end{proposition}

the isomorphism $a_F : A \xrightarrow{\sim} \Gr^W\!A$ respecte the filtration $W$, and trans-
form $F$ en $F_W$, and $\bar{F}$ en $d^{-1}(F_W)$: one has $a_{\bar{F}} = d^{-1} a_F$. The functor $\Gr^W$
admet therefore for inverse \`{a} gauche the functor $s : (B = \bigoplus B^{p,q}, d) \mapsto
(B, W, F_W, d^{-1}(F_W))$ equipped with the isomorphism fonctoriel $a_F : A \xrightarrow{\sim} s\Gr^W\!A$.

R\'{e}ciproquement, partons of a object bigradu\'{e} $B$, equipped with a automor-
phisme $d$, and d\'{e}finissons $W$, $F$, and $\bar{F}$ as ci-dessus. These filtrations
are oppos\'{e}es d\`{e}s that $\Gr^W(d) = 1$. For these filtrations, $B_F^{p,q} = B^{K(p,q)} \cap
d^{-1}(B^{I(p,q)})$. Si $d$ v\'{e}rifie \eqref{eq:1.1.1}, one has $d B^{p,q} \subset B^{I(p,q)}$, $d^{-1} B^{p,q} \subset B^{J(p,q)}$.
Puisque $(B^{p,q})$ and $(B_F^{p,q})$ are of the bigraduations of $B$, one has \'{e}galit\'{e}: $B_F^{p,q} =
B^{p,q}$. The m\^{e}me argument, appliqu\'{e} \`{a} $d^{-1}$ and \`{a} the bigrading $\bar{F}^{p,q} =
(B^{\bar{F}^{p,q}})$, shows that $d^{-1} B^{\bar{F}^{p,q}} = B^{p,q}$. One en d\'{e}duit that the isomorphism
\'{e}vident of $B$ with $\Gr^W\!B$ is a isomorphism of functors $\mathrm{Id} \xrightarrow{\sim} \Gr^W \circ s$.

\begin{remark}\label{rem:1.3}
Si $2$ is inversible, $d$ admet a unique racine carr\'{e}e $d^{1/2}$
v\'{e}rifiant again \eqref{eq:1.1.1}. Posons $a := d^{1/2} a_F = d^{-1/2} a_{\bar{F}}$. one has $a_{\bar{F}} = d^{1/2} a$
and $a_F = d^{-1/2} a$, of sorte that $a$ transforme $F$ (resp.~$\bar{F}$) en the filtration
$d^{1/2}(F_W)$ (resp.~$d^{-1/2}(\bar{F}_W)$) of $\Gr^W\!A$. the isomorphism $a$ is the m\^{e}me
for $A$ equipped with $W$, $F$, and $\bar{F}$ or of $W$, $\bar{F}$, and $F$.
\end{remark}

\subsection{}

These constructions are compatible to the tensor product: si $\otimes :
\A_1 \times \A_2 \to \A$ is a functor biadditif exact, one the \'{e}tend to the objects filtr\'{e}s
en posant
\begin{equation*}
F^n(A_1 \otimes A_2) = \sum_{a+b=n} F^a(A_1) \otimes F^b(A_2) \subset A_1 \otimes A_2.
\end{equation*}
The functor ``gradu\'{e} associ\'{e}'' is compatible \`{a} $\otimes$ and, for the objects bifiltr\'{e}s,
the isomorphism $\Gr^{W_1 \otimes W_2}(A_1 \otimes A_2) \xrightarrow{\sim} \Gr^{W_1}(A_1) \otimes \Gr^{W_2}(A_2)$ transforme the
filtration induite by $F_1 \otimes F_2$ en the $\otimes$ of celles induites by $F_1$ and $F_2$.

For the vector spaces over a field, that can se d\'{e}duire of this that a
bifiltration is toujours under-jacente \`{a} a bigrading. One en d\'{e}duit that si
$W_i$, $F_i$, and $\bar{F}_i$ are trois filtrations oppos\'{e}es on $A_i$ ($i=1,2$), then
the filtrations $W_1 \otimes W_2$, $F_1 \otimes F_2$, and $\bar{F}_1 \otimes \bar{F}_2$ of $A_1 \otimes A_2$ are again oppos\'{e}es.

The formation of the bigraduations $A_F^{p,q}$ and $A_{\bar{F}}^{p,q}$ of $A$ and $A_{W}^{p,q}$ of $\Gr^W\!(A)$, ainsi
that celle of the isomorphisms $a_F$, $a_{\bar{F}} : A \xrightarrow{\sim} \Gr^W\!(A)$ and of the automorphism
$d$ of $\Gr^W\!(A)$ is compatible \`{a} $\otimes$.

\subsection{}

Let $k$ a field commutative, $\A$ the cat\'{e}gorie of the spaces vecto-
riels (under-entendu: of dimension finite) on $k$ and $\mathscr{C}$ the cat\'{e}gorie of the vector spaces equipped with trois filtrations oppos\'{e}es $W$, $F$, and $\bar{F}$. C'is a $\otimes$-cat\'{e}gorie ab\'{e}lienne rigide, with $\End(1) = k$. It admet the functor fiber
$\omega_0$: ``space vector under-jacent'', therefore is tannakienne and neutre. One dispose also of the functor fiber $\Gr^W =: \omega_W$. The isomorphisms $a_F$ and $a_{\bar{F}}$ are
of the isomorphisms of functors fibers: $\omega_0 \xrightarrow{\sim} \omega_W$. Posons $G = \Aut(\omega_W)$.
the automorphism fonctoriel $d$ d\'{e}finit a point of the group pro-alg\'{e}brique $G$
on $k$.

\`{A} the bigrading fonctorielle of $\Gr^W(V)$ by the $V_W^{p,q}$ correspond a
homomorphisme
\begin{equation*}
\alpha : \Gm \times \Gm \longrightarrow G.
\end{equation*}
C'is a section of the projection
\begin{equation*}
\pr : G \to \Gm \times \Gm
\end{equation*}
corresponding \`{a} the inclusion in $\mathscr{C}$ of the cat\'{e}gorie of the vector spaces
bigradu\'{e}s (by $B \mapsto (B, W, F_W, \bar{F}_W)$). Je d\'{e}sire normaliser $\alpha$ of sorte
that $\alpha(\lambda, \mu) \cdot v = \lambda^p \mu^{-q} v$ for $v \in V_W^{p,q}$.

of apr\`{e}s the proposition~\ref{prop:1.2}, $G$ is also the group alg\'{e}brique of the automorphismes of the functor fiber ``space vector under-jacent'' on the cat\'{e}gorie of the
vector spaces bigradu\'{e}s equipped with a automorphism $d$ v\'{e}rifiant \eqref{eq:1.1.1}.
We shall the ``calculer'' lorsque $k$ is of caract\'{e}ristique $0$.

Let $\mathcal{L}$ the alg\`{e}bre of Lie on $k$, librement engendr\'{e}e by of the \'{e}l\'{e}ments
$D^{i,j}$ ($i, j < 0$), and munissons-the of the unique bigrading for which
$D^{i,j}$ is of bidegr\'{e} $(i, j)$. The $W_n(\mathcal{L})$ are of the id\'{e}to the of $\mathcal{L}$, and the quotients $\mathcal{L}/W_n(\mathcal{L})$ are of the alg\`{e}bres of Lie nilpotentes. We shall denote $U$ the
group pro-alg\'{e}brique projective limit of the groups alg\'{e}briques unipotents
$\exp(\mathcal{L}/W_n(\mathcal{L}))$ of alg\`{e}bre of Lie the $\mathcal{L}/W_n(\mathcal{L})$. Faisons agir the group $\Gm \times \Gm$
on $\mathcal{L}$, of sorte that $(\lambda, \mu)$ agisse on $D$ of degr\'{e} $(i, j)$ by multiplication
by $\lambda^i \mu^j$. This action respecte the filtration $W$, and fournit therefore a action
of $\Gm \times \Gm$ on $U$.

\begin{construction}\label{constr:1.6}
$G$ is the product semi-direct $(\Gm \times \Gm) \ltimes U$.
\end{construction}

Faire agir the product semi-direct on a space vector $V$ of dimension
finite revient \`{a} faire agir $\Gm \times \Gm$, and $U$, with a condition of compatibilit\'{e}
reliant these deux actions. the action of $\Gm \times \Gm$ correspond \`{a} a bigrading
(celle for which $(\lambda, \mu) \cdot v^{p,q} = \lambda^p \mu^{-q} v^{p,q}$). Celle of $U$ \`{a} a action
of $\mathcal{L}$, zero on a $W_n$ and nilpotente, i.e., \`{a} the donn\'{e}e of endomorphisms
$D^{i,j}$, for $i, j < 0$, such that all compos\'{e} of $D^{i,j}$ of weight sufficiently grand en
value absolute let zero. The compatibilit\'{e} is that $D^{i,j}$ let of bidegr\'{e} $(i, j)$,
and this propri\'{e}t\'{e} entra\^{i}ne the annulations requises ci-dessus. The donn\'{e}e of the
$D^{i,j}$ \'{e}quivaut \`{a} celle of $D := \sum D^{i,j}$: to the total, the action of the product semi-direct
on $V$ s'identifie \`{a} the donn\'{e}e of a bigrading, and of a endomorphism
$V$\'{e}rifiant $D V^{p,q} \subset \bigoplus_{i<0, j<0} V^{p+i,q+j}$. Posant $d = \exp(D)$, one the identifie again
\`{a} the donn\'{e}e of a bigrading, and of a automorphism $d$ v\'{e}rifiant \eqref{eq:1.1.1}.
This construction is compatible to the tensor product, of o\`{u} the isomorphism
voulu $(\Gm \times \Gm) \ltimes U \xrightarrow{\sim} G \xrightarrow{\sim} \Aut(\omega_W)$. This isomorphism is bien s\^{u}r
compatible \`{a} $\alpha$ and $\pr$ ci-dessus.

\section{Mixed Hodge structures}

Recall that a mixed Hodge structure r\'{e}it is a space vector
r\'{e}el (under-entendu: of dimension finite) equipped with a filtration croissante $W$
and of which the complexifi\'{e} $V_\C$ is equipped with a filtration d\'{e}croissante $F$. One
exige that on $V_\C$ the filtration $W_\C$---souvent not\'{e}e simplement $W$---d\'{e}duite
of $W$ by extension of scalars, the filtration $F$ and sa complex conjugu\'{e}e
$\bar{F}$ formant a syst\`{e}me of trois filtrations oppos\'{e}es.

It revient to the m\^{e}me of se donner a space vector r\'{e}el or a space vec-
toriel complex equipped with a structure r\'{e}it, i.e., of a conjugaison complex
(c'is a involution antilin\'{e}aire). Via this dictionnaire, it revient to the m\^{e}me of
se donner a mixed Hodge structure r\'{e}it, or a space vector com-
plexe equipped with trois filtrations oppos\'{e}es $W$, $F$, and $\bar{F}$, and of a conjugaison
complex which respecte $W$ and \'{e}change $F$ and $\bar{F}$. of apr\`{e}s the proposition~\ref{prop:1.2},
the functor $\Gr^W$ induces a \'{e}quivalence of the cat\'{e}gorie of these objects with
celle of the vector spaces complexes bigradu\'{e}s $V = \bigoplus V^{p,q}$, equipped with a
automorphism $d$ v\'{e}rifiant \eqref{eq:1.1.1} and of a conjugaison complex $c$ which
\'{e}change $V^{p,q}$ and $V^{q,p}$, and for which the complex conjugu\'{e} $\bar{d}$ of $d$ is
$d^{-1}$: $\bar{d} = c d c^{-1} = d^{-1}$.

The cat\'{e}gorie $\mathscr{H}$ of the mixed Hodge structures r\'{e}they is a $\otimes$-cat\'{e}gorie
ab\'{e}lienne rigide, with $\End(1) = \R$. It admet the functor fiber $\omega_0$: ``es-
pace vector under-jacent'', therefore is tannakienne and neutre. One dispose also
of the functor fiber $\Gr^W =: \omega_W$. With the notations of the remark~\ref{rem:1.3}, the
morphism fonctoriel
\begin{equation*}
a : V_\C \longrightarrow \Gr^W(V_\C)
\end{equation*}
respecte the structure r\'{e}it of the deux membres (car it d\'{e}pend sym\'{e}triquement
of $F$ and $\bar{F}$). It induces a isomorphism of functors fibers, again not\'{e}
$a : \omega_0 \xrightarrow{\sim} \omega_W$.

With the notations of the num\'{e}ro~1, faisons $k = \C$, and let $G_2$ the form r\'{e}it
of $G$ d\'{e}finite by the conjugaison complex suivante:
\begin{enumerate}[label=\alph*)]
\item on $\Gm \times \Gm$, c'is $(\lambda, \mu) \mapsto (\bar{\mu}, \bar{\lambda})$;
\item on $U$, c'is celle which is compatible \`{a} the conjugaison complex of $\mathcal{L}$
don\'{e}e by
\begin{equation*}
D^{i,j} \mapsto -D^{j,i}.
\end{equation*}
\end{enumerate}

As en 1.1, one sees that it revient to the m\^{e}me of faire agir $G_2$ on a
space vector r\'{e}el $V$, or of se donner
\begin{enumerate}[label=\alph*)]
\item a bigrading $V_\C = \bigoplus V^{p,q}$ of $V_\C$, such that $V^{p,q}$ and $V^{q,p}$ let
complexes conjugu\'{e}s the a of the other, and
\item a automorphism $d$ of $V_\C$, such that $\bar{d} = d^{-1}$, and that
\begin{equation*}
(d-1)V^{p,q} \subset \bigoplus_{\substack{i < p \\ j < q}} V^{i,j}.
\end{equation*}
\end{enumerate}

D\`{e}s lors,

\begin{proposition}\label{prop:2.1}
The functor fiber $\omega_W$ induces a \'{e}quivalence of $\mathscr{H}$ with
the cat\'{e}gorie of the repr\'{e}sentations of $G_2$: one has $\mathscr{H} = \Aut(\omega_W)$.
\end{proposition}

The functor inverse $\Phi$ is

$(V \text{, muni de l'action de } G_2) \mapsto$

$(V \text{, muni d'une bigraduation de } V_\C \text{, et de } d) \mapsto$

$(V \text{, muni de la filtration } W \text{ de complexifi\'{e}e la filtration } W \text{ attach\'{e}e \`{a} la}$
bigrading of $V_\C$, and of the filtration $F := d^{-1/2}(F_W)$ of son complexifi\'{e}$)$,

with for isomorphism $\Phi \circ \Gr^W \xrightarrow{\sim} \mathrm{Id}$ the isomorphism fonctoriel $a$.

\vspace{1cm}

\begin{flushright}
\textsc{Institute for Advanced Study, Princeton, New Jersey}
\end{flushright}

\vfill

\begin{center}
\rule{0.5\textwidth}{0.4pt}
\end{center}

\vspace{0.3cm}

{\small\noindent\textbf{1991 Mathematics Subject Classification.} Primary 14Fxx.}

\vspace{0.2cm}

{\small\noindent The pr\'{e}sent article reproduit with of the changements minimes, the appendice portant the m\^{e}me titre
to the preprint of E.~Cattani and A.~Kaplan: ``One the SL(2)-orbits in Hodge theory'' (IHES 1982).
A pr\'{e}sentation diff\'{e}rente of a part of the r\'{e}sultats is donn\'{e}e to the pages 470--473 of the article
of E.~Cattani, A.~Kaplan, and W.~Schmid: ``Degeneration of Hodge structures'' (Annals of Math.
\textbf{123} (1986), 457--535).}

\vspace{0.2cm}

{\small\noindent This paper is in final form and no version of it will be submitted for publication elsewhere.}

\vspace{0.2cm}

{\small\noindent \copyright 1994 American Mathematical Society\\
0082-0717/94 \$1.00 + \$.25 per page}

\end{document}
